\section{Conclusões}\label{sec:conclusoes}

Este trabalho apresenta um procedimento de emulação distribucional dependente do tempo para a análise da margem de durabilidade de elementos de concreto armado sob carbonatação. O encadeamento entre GLD, PCE e redes neurais permite representar distribuições condicionais e consultar sua evolução temporal. O estado limite de referência é a despassivação, e a comparação de limiares positivos expressa políticas de intervenção preventiva.

No problema de referência, as comparações locais em 50 e 100 anos apresentam divergências KL estimadas inferiores a $0{,}004$ para PCE e rede neural. A análise dos quantis evidencia, contudo, que uma pequena discrepância global pode coexistir com erro relativo elevado no quantil inferior quando este se aproxima de zero. A mediana do ganho de tempo na obtenção dos parâmetros por consulta à PCE é de aproximadamente $5{,}57\times10^5$, em relação ao procedimento que inclui amostragem, organização dos dados e ajuste local da GLD. Esse valor caracteriza a consulta ao modelo já treinado, com os limites de cronometragem descritos no Exemplo~1.

\falta{Concluir a avaliação da aplicação de durabilidade com os resultados dos Exemplos~2 e~3: acurácia da PCE e das redes, comportamento das caudas, verificação temporal, fração censurada, tempos até os limiares e custo computacional. Os resultados do problema de referência não substituem essa avaliação.}

A interpretação da vida remanescente depende das hipóteses de variabilidade latente, da fixação dos parâmetros de forma da GLD e do acoplamento temporal das trajetórias. Também é limitada pelo domínio de aplicação do modelo de carbonatação, pela resolução temporal e pela censura ao final do horizonte. A comparação de limiares descreve o prazo até uma ação preventiva; a avaliação de ganho de vida após reparo requer um modelo da intervenção e da degradação posterior.

Como continuidade, propõem-se a caracterização das incertezas com dados de inspeção, a verificação do emulador em novos pontos e idades, a avaliação da dependência temporal e a incorporação de modelos de reparo e custos. Essas etapas permitem avançar da comparação de políticas para a seleção de ações de manutenção em um ativo específico.
